\section{Threats to validity}
\label{sec:threats}
% Built from docs/plan/two-track-plan.md section 2.6 (threats and mitigations) and section 9
% (reviewer and scope critiques), plus threats that follow from the code and the TO-OCP
% formulation (sections 3 and 4 of this draft). Each paragraph states the threat, what
% the design does about it, and what remains.

\paragraph{The labels are local optima} IPOPT returns local optima, so every label is a
lower bound (\cref{sec:answer-lb}). A lower bound is harmless for ranking only if the
solver's shortfall is unrelated to what the scores respond to. If the solver fails more
often on, say, hairpin-heavy levels, measured regret is biased low exactly there. The
uniform three-start budget, including a start from the agent's own best rollout, the
200-level repeatability study and the check against observed agent returns bound the
shortfall; we report label quality per level set and per track-type stratum
(\cref{sec:results-labels}).
% src: docs/plan/two-track-plan.md sections 2.2 and 2.6 (row "Labels are local optima")

\paragraph{Labels depend on the agent} One of the three warm starts is the final
checkpoint's best rollout, so how tight a label is depends on how well that agent drives
the level. Label error is then correlated with agent-derived quantities, including the
agent returns that MaxMC's $\Rmax$ is built from: on levels where the final agent is
strong, its rollout lifts the label, and a score that tracks the agent's own returns can
look better aligned with regret than it is. The uniform budget keeps the effect the same
for every level, but does not remove it. We pre-register a report of the label gain from
the agent start, $\Jstar$ minus the best of the two agent-free starts, by level set and
track-type stratum, and repeat the headline endpoints with labels from the agent-free
starts only.
\TODO{add both to \code{protocol-v1}}

\paragraph{Label noise late in training} As the agent improves, regret shrinks towards
the size of the label noise and any correlation is attenuated. A late-phase degradation
is claimed only where median regret is at least three times the label standard deviation,
correlations are also reported after reliability correction, and the E4 trigger requires
both. If late-phase labels turn out too noisy, the paper falls back to early- and
mid-phase claims and the buffer audit.
% src: docs/plan/two-track-plan.md sections 2.2 ("Label noise"), 2.5 ("E4 trigger"), 7

\paragraph{Regret no policy in the class can remove} Part of $\Reg$ is out of reach for
our agent: it does not observe the hidden grip draw ($\pm2\%$), sees at most 500\,m ahead,
and acts on noisy observations. This is regret in the sense UED defines it, but a
regret-maximising sampler can fixate on exactly such levels~\citep{beukman2024remidi}. The
design measures it rather than removing it: E0-short changes how much the agent can see,
the oracle arm in C is compared with DR on held-out levels, and the E4-switch follow-up
has a control arm scored against $\Jhund$.
% src: src/apexzero/env/adapters.py (mu_jitter 0.02); src/apexzero/env/racing.py (preview_offsets);
%      docs/plan/two-track-plan.md section 2.5 (row B follow-up)
\TODO{estimate the share of regret due to the hidden grip by labelling a subset of levels
twice, with the drawn and with the nominal grip}

\paragraph{The finite window} The label maximises return over exactly \num{1200} steps,
while PPO bootstraps the window end and learns a time-unaware policy. Two effects follow:
the label omits the tail $\gamma^{T}V(x_T)$, and the optimiser may arrive at the window
end in a state no time-unaware policy would choose, such as too fast for a corner just
beyond it. Both are weighted by $\gamma^{1200} \approx 0.0024$.
% src: docs/paper/sources-jaxued.md section 4 item 3 (gamma^1200)
\RESULT{on a subset of levels, the change in $\Jstar$ when the window is extended by 10\,s
and the extra steps are dropped from the objective, as a bound on both effects}

\paragraph{An audit of a sampler nobody runs} Fresh scores on DR checkpoints are not what a
curriculum sees: a curriculum's buffer holds old scores from earlier policies, chosen by
an admission rule. C therefore runs real \PLRp{} and SFL, and in-buffer fidelity and
admission AUROC are co-primary. Our samplers are adapted to windows that span chunks
(\cref{tab:deviations}); E5 runs the reference code unmodified.
% src: docs/plan/two-track-plan.md section 2.6 (row "It audits a sampler nobody runs")

\paragraph{A privileged critic flatters value-based scores} E0's critic sees more than the
actor, which may make PVL and MaxMC look better than in standard UED. E0-sym, with a
symmetric critic, is co-primary, and the size of the drop is reported.
% src: docs/plan/two-track-plan.md section 2.6
\TODO{the asymmetric critic is not implemented yet; until it is, all runs are E0-sym}

\paragraph{Training schedule moves with phase} Learning rate and entropy coefficient decay
during training, so a change in fidelity with phase could be a change of schedule. E0-const
holds both fixed, schedule variables enter as covariates, and deterministic-policy windows
are reported separately.
% src: docs/plan/two-track-plan.md section 2.6

\paragraph{Discount weighting and mixed windows} With $\gamma = 0.995$ the first seconds of
a window dominate its return, and windows that crash mix with clean ones. Levels start on
a straight; strata and score variants are discount-weighted; the undiscounted
maximum-distance target removes the weighting; and the hurdle split separates crash
probability from regret given completion.
% src: docs/plan/two-track-plan.md section 2.6

\paragraph{One bespoke domain with dense reward} Racing is one domain, and its reward is
dense. The maze replication uses JaxUED's own code and sparse reward; the critic and
preview manipulations vary the racing agent itself. A sparse-reward racing arm was
considered and rejected for this version, because it would need PPO re-tuned on an
untested reward; claims about reward density are therefore stated as conjecture.
% src: docs/plan/two-track-plan.md sections 2.6 and 9 ("One bespoke domain. Partly accepted")

\paragraph{Forking paths} Six headline estimators are fixed in advance, with Holm
correction within the co-primary family; \code{audit-v1} and \code{protocol-v1} are tagged
before any full-scale result; and the claim ledger is written by the figure scripts, not by
hand.
% src: docs/plan/two-track-plan.md section 2.6

\paragraph{Numerical artefacts} A label is wrong if the optimiser exploits a difference
between its model and the simulator. Both evaluate the same $C^2$ spline and the same step
map (\cref{tab:gates}), every label is replayed in float64, and a float32-against-float64
error bar covers the fact that agents run in single precision.
% src: docs/plan/two-track-plan.md section 2.6 (row "Numerical artefacts")
\TODO{measure the float32-vs-float64 error bar}

\paragraph{The optimiser finds what the model allows} An optimum is only as sensible as the
model: an early minimum-time solution drifted far beyond the rear tyre's peak slip, on
the soft post-peak tyre curve, which we then closed with a slip limit applied identically
in the environment and the optimiser (\cref{sec:env-mdp}). An agent could find the same
exploits, and then they are part of both sides of $\Reg$, but we want neither side to
depend on them. Labels are computed only after the simulator is frozen as \code{sim-v1},
and any exploit found later is reported, not silently fixed.
% src: docs/STATUS.md ("Things found along the way"); docs/plan/two-track-plan.md section 4 (K5, K2b)

\paragraph{Realism is not the point} The vehicle uses literature values where they exist
and assumptions elsewhere (\cref{tab:params}); it is an F1-style car, not a calibrated
one. The audit does not need realism. It needs the agent's and the optimiser's problems
to be identical, which the parity, integrator and replay gates check. Realism only
supports the weaker claim that the domain exercises racing-like control, which the
terminal-speed, braking and skidpad gates support.

\paragraph{Lockbox hygiene} Twelve real circuits are a lockbox, held out of this paper
for later evaluation, and twelve disjoint ones are for development. No audit or
curriculum experiment uses a real circuit. The generator is fitted on the development
circuits only, and the development checks (the specialist, the minimum-time replay gate
of \cref{tab:gates} and the solver spike) ran on development circuits only. No number in
the paper is computed on a lockbox circuit. Two departures from a clean lockbox happened
before the simulator was frozen, and we disclose both. First, minimum-time optima were
once solved on every TUMFTM circuit, the lockbox included, against a plan that allows
optimal control only on development circuits until the generator is frozen; the replay
gate reported here was recomputed on the development circuits alone, and none of those
lockbox results is used. Second, an earlier minimum-time optimum on a lockbox circuit
revealed the slip exploit above, and the $10^\circ$ slip limit was introduced because of
it. A lockbox circuit therefore shaped one simulator design choice. The exploit is a
property of the tyre curve, which has the same post-peak shape on every track, but the
decision to close it came from a lockbox result, so any later lockbox evaluation must
report it.
% src: docs/STATUS.md (K2a row; "Things found along the way", first item);
%      src/apexzero/tracks/corners.py (R_DEV, R_TEST);
%      docs/plan/two-track-plan.md section 4 ("OCPs run on R-dev only until generator-v1 is tagged",
%      "The paper never touches R-test"); R-dev recomputation: see the src comment of tab:gates
\VERIFY{the docs record the exploit and the fix (docs/STATUS.md) but not how the value
$10^\circ$ was chosen; if it was tuned against the lockbox optimum, say so here}
Separately, the public telemetry of \cref{app:snapping} covers qualifying sessions at
real circuits, lockbox circuits among them. It is human data only: no agent, optimum,
quasi-steady-state profile or simulator quantity is computed for any circuit there, and
no claim depends on it.
